Quatre càrregues estan situades en els vèrtexs d'un quadrat de 4,00~m de costat, tal com s'indica en la figura. Els valors de les càrregues són $Q_1 = 1,00\ \mathrm{\mu C}$, $Q_2 = -2,00\ \mathrm{\mu C}$, $Q_3 = 2,00\ \mathrm{\mu C}$ i $Q_4 = -1,00\ \mathrm{\mu C}$. El punt C és a la intersecció de les dues diagonals. El punt A està situat a la meitat del segment que va des de la càrrega $Q_1$ fins a la càrrega $Q_2$.

\figura[0.4]{fig1.png}

\begin{apartat}{1}
Representeu i calculeu el vector camp elèctric en el punt C.
\end{apartat}

\begin{apartat}{1}
Calculeu la diferència de potencial entre els punts A i C.
\end{apartat}

\blocetiquetat{Dada:}{$k = \dfrac{1}{4\pi\varepsilon_0} = 8,99\times 10^{9}\ \mathrm{N\,m^2\,C^{-2}}$.}
